\section*{前言}
\addcontentsline{toc}{section}{前言}

\subsection*{这本书想回答什么}

1989 年，Witten 在《Communications in Mathematical Physics》上发表了一篇四十页的论文，
标题是《Quantum Field Theory and the Jones Polynomial》。论文的中心论断只有一句话：
Jones 多项式是一个三维量子场论的配分函数。这句话之所以重要，不是因为它给了多项式一个新的算法，
而是因为它把“纽结”从一个二维组合对象变成了三维拓扑的一部分，并由此把纽结理论、量子场论与
共形场论放进同一个框架。

本讲义试图把这句话讲成一个连贯的故事，而不是逐条注释一篇论文。我们只假定读者学过量子力学
和一学期的场论或微分几何；共形场论、几何量子化、仿射李代数这些工具，都会在用到的时候现讲。
读完本书，读者应当能独立地把这条链复述清楚：
\[
\begin{aligned}
&\text{三维 Chern--Simons 理论}
\;\longrightarrow\;
\text{曲面上的量子 Hilbert 空间}
\;\longrightarrow\;
\text{模空间上的全纯截面}\\
&\qquad\longrightarrow\;
\text{共形块}
\;\longrightarrow\;
\text{融合规则与 Verlinde 公式}.
\end{aligned}
\]

\subsection*{读者指南}

这本讲义按“一条主线、两个侧面”组织。主线是：Chern--Simons 理论在一个曲面上量子化，
得到的态空间恰好是 WZW 模型的共形块。两个侧面是：三维侧的 Wilson 线与手术计算，
二维侧的融合规则与 Verlinde 公式。读完八讲，读者应当能自己把这条主线复述一遍，
并知道每一步在什么条件下成立。

前置知识只有三样：量子力学（知道 Hilbert 空间与算符）、一门场论课（知道路径积分与规范场），
以及微分几何的基本语言（流形、丛、联络）。辛几何与复几何的概念在用到时现讲，
不需要提前准备。如果读者只想知道结论，读前言、第一讲与第八讲即可；
如果要动手算，第七讲是入口。

建议的读法是“两遍”：第一遍只读每讲开头的动机与结尾的小结，跳过所有带编号的推导；
第二遍再逐条推公式。书中标注 \Wp{n} 的页码是原文的印刷页，仅供核对，
不影响阅读顺序。

\subsection*{路线图}

全书分八讲。第一讲从 Jones 多项式的组合定义出发，说明它为什么需要一个三维的解释。
第二讲引进 Chern--Simons 作用量，并讲清楚它的两个反直觉性质：不依赖度规，以及耦合常数必须取整数。
第三讲把时空沿曲面切开，读出相空间：曲面上平坦联络的模空间。第四讲把相空间量子化，
这是全书技术上的核心，几何量子化的预量子化、极化、全纯截面三步都在这里。
第五讲把纽结放进来，讲 Wilson 线如何变成穿刺点上的一个约束，以及这个约束为什么不能经典地解。
第六讲转向二维侧，讲 WZW 模型、共形块与融合规则。
第七讲回到三维，用手术把配分函数算出来，并推出 Verlinde 公式。
第八讲讨论两个常见的误读——把本文当作全息对偶的雏形，或当作范畴化的预言——并给出仍然开放的问题。

\subsection*{约定}

\begin{itemize}[leftmargin=*]
  \item $M$ 表示紧致、定向、无边三维流形，$\Sigma$ 表示其中的曲面。$G$ 是紧致单 Lie 群，
        $d=\dim G$，$h^\vee$ 是它的对偶 Coxeter 数（$SU(N)$ 时 $h^\vee=N$）。
  \item 本文的 Chern--Simons 作用量取
        $L=\frac{k}{4\pi}\int_M\Tr(A\wedge\dd A+\frac23A\wedge A\wedge A)$，
        其中 $\Tr$ 取定义表示、长根平方为 $2$。$k$ 是整数。
  \item 引用原文时用边注 \Wp{370}，表示 Witten 论文的印刷页 370。
        教材页码在脚注里给出。
  \item 我们区分三种陈述：原文给出的构造、我们在正文里补出的推导、以及仍然只能算猜想的部分。
        正文中凡属后两类，都会明确写出“可验证的后果”，读者可以自己核对。
\end{itemize}
